\scanpage{28}
\mspage{26}
% archive sheet number: 27
\deleted{Pour prouver} \inserted{Dire} que les produits finis indexés par l'ens. fini $I$ existent dans la catégorie $\mathcal A=\mathbb D(e)$ (disons) équivaut à dire que le foncteur diagonal
\[
\mathcal A\xrightarrow{\Delta}\mathcal A^I
\]
admet un adjoint à gauche. Or si
\[
f:I\longrightarrow e
\]
est la projection de $I$ (\deleted{\unsure{con}}considéré comme ens. ord. discret) sur \deleted{\unsure{le}} l'\unsure{en|ens} ordonné ponctuel $e$, le foncteur $\Delta$ s'identifie à $f^*:\mathbb D(e)\longrightarrow\mathbb D(I)$ (grâce à l'axiome $\unsure{\mathrm{D\bar er}|\mathrm{D\acute er}|\mathrm{Der}}_2$ sur le \unsure{comportement} des dérivateurs pour une somme \unsure{directe|discrète} de cat. d'indices).
Donc l'existence des produits revient à celle de l'adj. à gauche $f_*$ de $f^*$, \inserted{OK pour dérivateur g.} Et un hom. exact à gauche de dérivateurs à gauche, comme il commute aux $f_*$, commute aux produits finis. \unsure{Item|Itou} pour les sommes finies, et les dérivateurs gauches.

L'exemple précédent utilisait le fait que le dérivateur envisagé n'est pas pointé. Mais Même pour un dérivateur pointé \inserted{$\mathbb D$}, un morphisme $f^{\mathbb D}_*:\mathbb D_I\longrightarrow\mathbb D_{I'}$ n'est pas forcément exact à droite. Si on prend le dérivateur associé à une catégorie $\mathcal M$ quelc. (avec $\Sigma=\varnothing$),
\[
\mathbb D_{\mathcal M}(I)=\underline{\mathrm{Hom}}(I^\circ,\mathcal M)\overset{\text{déf}\unsure{\text{n}|\,}}{=}\mathcal M(I)
\]
alors le dérivateur induit
\[
\begin{aligned}
(\mathbb D_{\mathcal M})_I&=\big\{J\longmapsto\mathbb D_{\mathcal M}(I\times J)=\underline{\mathrm{Hom}}((I\times J)^\circ,\mathcal M))\\
&\hspace{6em}\simeq\underline{\mathrm{Hom}}(J^\circ,\underline{\mathrm{Hom}}(I^\circ,\mathcal M))
\end{aligned}
\]
